\section*{Chapitre 6 \textemdash\ Droites et Cercles}
\addcontentsline{toc}{section}{Chapitre 6 \textemdash\ Droites et Cercles}

\subsection*{Exercices}
\addcontentsline{toc}{subsection}{Exercices}

%────────────────────────────────────────────────────────────────
\solutiontitle{6}{1}{\'Equation d'une droite}

\begin{enumerate}
  \item $m=\frac{9-3}{5-2}=2$~; $p=3-4=-1$~: $y=2x-1$.
  \item $y-4=-2(x+1)\Rightarrow y=-2x+2$.
  \item Pente $2$, par $Q(3,0)$~: $y=2(x-3)=2x-6$.
  \item $x_R=x_S=1$~: droite verticale $x=1$.
  \item Perp. \`a pente $3$~: pente $-\frac{1}{3}$. Par $A(0,-1)$~: $y=-\frac{x}{3}-1$.
  \item $m=\frac{-2-4}{4-(-2)}=-1$~; $p=4-(-1)(-2)=2$~: $y=-x+2$.
\end{enumerate}

%────────────────────────────────────────────────────────────────
\solutiontitle{6}{2}{Droite~: forme cart\'esienne}

\begin{enumerate}
  \item $y=2x+3$~; $m=2$, $p=3$.
  \item $y=-\frac{x}{3}+2$~; $m=-\frac{1}{3}$, $p=2$.
  \item $y=2x+4$~; $m=2$, $p=4$.
  \item $y=-5x$~; $m=-5$, $p=0$.
  \item $y=\frac{3}{4}x+3$~; $m=\frac{3}{4}$, $p=3$.
  \item Droite verticale $x=-2$~: pas de forme r\'eduite.
\end{enumerate}

%────────────────────────────────────────────────────────────────
\solutiontitle{6}{3}{\'Equation d'un cercle}

\begin{enumerate}
  \item $(x-1)^2+(y+2)^2=9$.
  \item $x^2+y^2=25$.
  \item $A(4,-2)$~: $(4-1)^2+(-2+2)^2=9+0=9$.\checkmark\quad Sur le cercle.
        $B(0,0)$~: $(0-1)^2+(0+2)^2=1+4=5\ne9$. Pas sur le cercle.
  \item $C(3,4)$~: $9+16=25$.\checkmark\quad Sur le cercle.
\end{enumerate}

%────────────────────────────────────────────────────────────────
\solutiontitle{6}{4}{Identifier centre et rayon}

\begin{enumerate}
  \item $(x-3)^2+(y+2)^2=9+4+3=16$. Centre $(3,-2)$, rayon $4$.
  \item $(x+1)^2+y^2=1+8=9$. Centre $(-1,0)$, rayon $3$.
  \item Centre $(3,0)$, rayon $4$.
  \item $(x-2)^2+(y-3)^2=4+9-4=9$. Centre $(2,3)$, rayon $3$.
\end{enumerate}

%────────────────────────────────────────────────────────────────
\solutiontitle{6}{5}{Distance d'un point \`a une droite}

\begin{enumerate}
  \item $d=\frac{|4(-2)-3(3)+5|}{\sqrt{16+9}}=\frac{|-8-9+5|}{5}=\frac{12}{5}$.
  \item $d=\frac{|3+4-5|}{5}=\frac{2}{5}$.
  \item $d=\frac{|0-0+26|}{13}=2$.
  \item $d=\frac{|2-1-1|}{\sqrt{2}}=0$~: le point est sur la droite.
  \item Forme cart\'esienne~: $2x-y-1=0$. $d=\frac{|6-4-1|}{\sqrt{5}}=\frac{1}{\sqrt{5}}=\frac{\sqrt{5}}{5}$.
  \item Droite verticale $x=4$~: $d=|3-4|=1$ (distance horizontale).
\end{enumerate}

%────────────────────────────────────────────────────────────────
\solutiontitle{6}{6}{Parall\`ele et perpendiculaire}

\begin{enumerate}
  \item Pente $3$, par $A(2,1)$~: $y=3x-5$.
  \item Pente perp. $-\frac{1}{2}$, par $B(-1,3)$~: $y=-\frac{x}{2}+\frac{5}{2}$.
  \item Droite $2x-y+1=0$~: pente $2$, perp. pente $-\frac{1}{2}$.
        Par $C(0,4)$~: $y=-\frac{x}{2}+4$.
  \item Pentes~: $m_1=3$, $m_2=-\frac{1}{3}$ (de $3y=-x+5$).
        $3\times(-\frac{1}{3})=-1$.\checkmark\quad Perpendiculaires.
\end{enumerate}

%────────────────────────────────────────────────────────────────
\solutiontitle{6}{7}{Intersection de deux droites}

\begin{enumerate}
  \item $2x-1=-x+5\Rightarrow3x=6$~: $x=2$, $y=3$. Point $(2,3)$.
  \item $x+3=3x-1\Rightarrow x=2$, $y=5$. Point $(2,5)$.
  \item $2x-y=4$ et $x+3y=9$~: $7x=21$, $x=3$, $y=2$. Point $(3,2)$.
  \item $y=4$~; $x-8+6=0\Rightarrow x=2$. Point $(2,4)$.
\end{enumerate}

%────────────────────────────────────────────────────────────────
\solutiontitle{6}{8}{Position relative droite/cercle}

\begin{enumerate}
  \item $\Omega(1,-2)$, $r=5$. $d=\frac{|3-4(-2)+2|}{5}=\frac{|3+8+2|}{5}=\frac{13}{5}=2{,}6<5$.
        \textbf{S\'ecante.}
  \item $\Omega(0,0)$, $r=3$, droite $x-y+4=0$.
        $d=\frac{|4|}{\sqrt{2}}=2\sqrt{2}\approx2{,}83<3$. \textbf{S\'ecante.}
  \item $\Omega(2,1)$, $r=2$, droite $y-3=0$.
        $d=|1-3|=2=r$. \textbf{Tangente.}
  \item $\Omega(0,0)$, $r=5$, droite $4x+3y-25=0$.
        $d=\frac{|{-25}|}{5}=5=r$. \textbf{Tangente.}
\end{enumerate}

%────────────────────────────────────────────────────────────────
\solutiontitle{6}{9}{M\'ediatrice d'un segment}

\begin{enumerate}
  \item Milieu $(3,2)$. $\vv{AB}=(4,-2)$~: perp. de pente $2$.
        $y-2=2(x-3)\Rightarrow y=2x-4$.
  \item Milieu $(2,2)$. $\vv{AB}=(4,4)$~: perp. de pente $-1$.
        $y-2=-(x-2)\Rightarrow y=-x+4$.
  \item Milieu $(1,-1)$. $\vv{AB}=(6,-4)$~: pente $-2/3$, perp. $3/2$.
        $y+1=\frac{3}{2}(x-1)\Rightarrow y=\frac{3x}{2}-\frac{5}{2}$.
  \item Milieu $(2,2)$. $\vv{AB}=(0,-6)$~: droite verticale, perp. horizontale $y=2$.
\end{enumerate}

%────────────────────────────────────────────────────────────────
\solutiontitle{6}{10}{Droites sp\'eciales d'un triangle}

$A(0,6)$, $B(-3,0)$, $C(3,0)$.
\begin{enumerate}
  \item $AB$~: pente $\frac{6}{3}=2$~; $y=2x+6$.
        $AC$~: pente $\frac{-6}{3}=-2$~; $y=-2x+6$.
        $BC$~: $y=0$ (axe des $x$).
  \item M\'ediatrice de $BC$~: milieu $(0,0)$, perp. \`a $y=0$~: $x=0$.
        M\'ediatrice de $AB$~: milieu $(-3/2,3)$, perp. \`a pente $2$ (pente $-1/2$)~:
        $y-3=-\frac{1}{2}(x+\frac{3}{2})\Rightarrow y=-\frac{x}{2}+\frac{9}{4}$.
        M\'ediatrice de $AC$~: milieu $(3/2,3)$, perp. \`a pente $-2$ (pente $1/2$)~:
        $y=\frac{x}{2}+\frac{9}{4}$.
  \item En $x=0$~: les deux m\'ediatrices obliques donnent $y=9/4$.
        Circocentre~: $O'=(0,9/4)$. V\'erification~:
        $O'A=|6-9/4|=15/4$~; $O'B=\sqrt{9+(9/4)^2}=\sqrt{9+81/16}=\frac{\sqrt{225}}{4}=\frac{15}{4}$.\checkmark
\end{enumerate}

%────────────────────────────────────────────────────────────────
\solutiontitle{6}{11}{Hauteur d'un triangle}

$A(1,5)$, $B(0,0)$, $C(4,2)$.
\begin{enumerate}
  \item $\vv{BC}=(4,2)$~: pente $\frac{1}{2}$. \'Equation de $BC$~: $y=\frac{x}{2}$.
  \item Perp. de pente $-2$ passant par $A(1,5)$~: $y-5=-2(x-1)\Rightarrow y=-2x+7$.
  \item Intersection~: $\frac{x}{2}=-2x+7\Rightarrow\frac{5x}{2}=7\Rightarrow x=\frac{14}{5}$~;
        $y=\frac{7}{5}$. $H=\bigl(\frac{14}{5},\frac{7}{5}\bigr)$.
\end{enumerate}

%────────────────────────────────────────────────────────────────
\solutiontitle{6}{12}{Tangente \`a un cercle}

\begin{enumerate}
  \item $A(3,4)$ sur $x^2+y^2=25$~: tangente $3x+4y=25$.
  \item Cercle $(x-1)^2+(y+2)^2=10$, $B(2,1)$.
        V\'erif.~: $(2-1)^2+(1+2)^2=1+9=10$.\checkmark
        Vecteur $\Omega B=(1,3)$~: tangente perp. (pente $-1/3$) passant par $B$~:
        $y-1=-\frac{1}{3}(x-2)\Rightarrow x+3y-5=0$.
  \item Pour le cercle unit\'e~: vecteur $\vv{\Omega A}=(3,4)$, tangente $3x+4y=25$.
        $\vv{\Omega A}\cdot(1,0)=3$~; tangente a vecteur directeur $(-4,3)$.
        $(3,4)\cdot(-4,3)=-12+12=0$.\checkmark Perpendiculaires.
\end{enumerate}

%────────────────────────────────────────────────────────────────
\solutiontitle{6}{13}{Cercle passant par trois points}

$x^2+y^2+Dx+Ey+F=0$. Substitution~:
$A(0,0)$~: $F=0$.
$B(4,0)$~: $16+4D+F=0\Rightarrow D=-4$.
$C(0,2)$~: $4+2E+F=0\Rightarrow E=-2$.
\'Equation~: $x^2+y^2-4x-2y=0\Rightarrow(x-2)^2+(y-1)^2=5$.
Centre $(2,1)$, rayon $\sqrt{5}$.

%────────────────────────────────────────────────────────────────
\solutiontitle{6}{14}{Lire une figure}

$\mathcal{C}$~: $(x-2)^2+(y-1)^2=9$~; $(\Delta)$~: $y=x-1$.

Forme cart\'esienne de $\Delta$~: $x-y-1=0$.
$d=\frac{|2-1-1|}{\sqrt{2}}=0<3$.
\textbf{S\'ecante.}

%────────────────────────────────────────────────────────────────
\solutiontitle{6}{15}{Droites et cercle~: intersection}

$x^2+y^2=25$, $y=x+1$.
\begin{enumerate}
  \item $x^2+(x+1)^2=25\Rightarrow2x^2+2x-24=0\Rightarrow x^2+x-12=0$.
  \item $\Delta=1+48=49$~; $x=\frac{-1\pm7}{2}$~: $x_1=3$, $x_2=-4$.
  \item $P_1=(3,4)$~; $P_2=(-4,-3)$.
  \item V\'erif.~: $9+16=25$.\checkmark\quad $16+9=25$.\checkmark
\end{enumerate}

%────────────────────────────────────────────────────────────────
\solutiontitle{6}{16}{Axe radical de deux cercles}

$\mathcal{C}_1$~: $x^2+y^2-4x+2y=0$~; $\mathcal{C}_2$~: $x^2+y^2+2x-6y+6=0$.
\begin{enumerate}
  \item Axe radical~: $(D_1-D_2)x+(E_1-E_2)y+(F_1-F_2)=0$.
        $(-4-2)x+(2-(-6))y+(0-6)=0\Rightarrow-6x+8y-6=0\Rightarrow3x-4y+3=0$.
  \item Si $A$ est une intersection~: $A\in\mathcal{C}_1$ et $A\in\mathcal{C}_2$,
        donc $\pi_1(A)=0=\pi_2(A)$~; $A$ est sur l'axe radical.\hfill$\square$
  \item $\mathcal{C}_1$~: centre $(2,-1)$, rayon $\sqrt{5}$.
        $\mathcal{C}_2$~: centre $(-1,3)$, rayon $2$.
        $d=\sqrt{9+16}=5>\sqrt{5}+2\approx4{,}24$~?
        $\sqrt{5}\approx2{,}24$~; $\sqrt{5}+2\approx4{,}24<5$~: les cercles sont ext\'erieurs.
        Pas d'intersection r\'eelle.
\end{enumerate}

%────────────────────────────────────────────────────────────────
\solutiontitle{6}{17}{Lieu g\'eom\'etrique}

\begin{enumerate}
  \item $MA=MB$~: m\'ediatrice de $[AB]$~: $x=2$ (droite verticale passant par le milieu $(2,0)$ de $[AB]$).
  \item $MA^2+MB^2=10$~: $(x-1)^2+y^2+(x+1)^2+y^2=10\Rightarrow2x^2+2y^2+2=10\Rightarrow x^2+y^2=4$.
        Cercle de centre $O$ et rayon $2$.
  \item $MA^2=4MB^2$~: $x^2+y^2=4[(x-3)^2+y^2]$.
        $x^2+y^2=4x^2-24x+36+4y^2\Rightarrow3x^2+3y^2-24x+36=0\Rightarrow(x-4)^2+y^2=\frac{16}{3}$... 
        \emph{Corr.~:} $3x^2-24x+3y^2+36=0\Rightarrow x^2-8x+y^2+12=0\Rightarrow(x-4)^2+y^2=4$.
        Cercle de centre $(4,0)$ et rayon $2$ (cercle d'Apollonius).
  \item $\vv{MA}\cdot\vv{MB}=0$~: $M$ est sur le cercle de diam\`etre $[AB]$.
        Centre $O=(0,0)$, rayon $2$~: $x^2+y^2=4$.
\end{enumerate}

%────────────────────────────────────────────────────────────────
\solutiontitle{6}{18}{Droite et cercle~: param\`etre}

$(\Delta_k)$~: $y=kx+1$~; $\mathcal{C}$~: $x^2+y^2=4$.
$d(O,\Delta_k)=\frac{|1|}{\sqrt{k^2+1}}$.
\begin{enumerate}
  \item S\'ecante~: $d<2\Rightarrow\frac{1}{\sqrt{k^2+1}}<2\Rightarrow k^2+1>\frac{1}{4}$~: toujours vrai.
        Mais $d<r$~: $\frac{1}{\sqrt{k^2+1}}<2\Rightarrow1<4(k^2+1)$~: toujours vrai.
        Donc la droite est s\'ecante pour tout $k\in\R$.
        \emph{Remarque~: la droite $y=1$ (k=0) est \`a distance $1<2$ du centre.}
  \item Tangente~: $d=2\Rightarrow\frac{1}{\sqrt{k^2+1}}=2\Rightarrow k^2+1=\frac{1}{4}$~: impossible ($k^2\geq0$).
        \textbf{Aucune tangente} de la forme $y=kx+1$ au cercle de rayon $2$.
  \item $k=0$~: $y=1$. $x^2+1=4\Rightarrow x=\pm\sqrt{3}$. Points~: $(\sqrt{3},1)$ et $(-\sqrt{3},1)$.
\end{enumerate}

%────────────────────────────────────────────────────────────────
\solutiontitle{6}{19}{Droite m\'ediatrice et circocentre}

$A(1,4)$, $B(-2,-2)$, $C(5,-2)$.
\begin{enumerate}
  \item M\'ediatrice de $[AB]$~: milieu $(-1/2,1)$, pente $AB=\frac{6}{3}=2$, perp. $-1/2$.
        $y-1=-\frac{1}{2}(x+\frac{1}{2})\Rightarrow2y=-x+\frac{3}{2}\Rightarrow x+2y=\frac{3}{2}$.
        M\'ediatrice de $[BC]$~: milieu $(3/2,-2)$, pente $BC=0$, perp. verticale $x=3/2$.
  \item En $x=3/2$~: $\frac{3}{2}+2y=\frac{3}{2}\Rightarrow y=0$. Circocentre $O'=(\frac{3}{2},0)$.
  \item $O'A=\sqrt{(\frac{1}{2})^2+16}=\sqrt{\frac{65}{4}}$~;
        $O'B=\sqrt{(\frac{7}{2})^2+4}=\sqrt{\frac{65}{4}}$~;
        $O'C=\sqrt{(\frac{7}{2})^2+4}=\sqrt{\frac{65}{4}}$.\checkmark
  \item Cercle~: $(x-\frac{3}{2})^2+y^2=\frac{65}{4}$.
\end{enumerate}

%────────────────────────────────────────────────────────────────
\solutiontitle{6}{20}{Hauteurs et orthocentre}

$A(0,4)$, $B(-2,0)$, $C(2,0)$.
\begin{enumerate}
  \item $BC$~: $y=0$. Hauteur de $A$~: perp. horizontale $y=0$ passant par $A$~: $x=0$ (verticale).
        $AC$~: pente $\frac{-4}{2}=-2$. Hauteur de $B$~: pente $1/2$, par $B(-2,0)$~:
        $y=\frac{1}{2}(x+2)=\frac{x}{2}+1$.
  \item Intersection~: $x=0$~; $y=1$. Orthocentre $H=(0,1)$.
  \item Hauteur de $C$~: $AB$ pente $\frac{4}{2}=2$, perp. $-1/2$, par $C(2,0)$~:
        $y=-\frac{1}{2}(x-2)=-\frac{x}{2}+1$. En $x=0$~: $y=1$.\checkmark
\end{enumerate}

%────────────────────────────────────────────────────────────────
\solutiontitle{6}{21}{Position relative de deux cercles}

\begin{enumerate}
  \item $d=5$, $r_1+r_2=4$, $|r_1-r_2|=2$. $d>r_1+r_2$~: \textbf{ext\'erieurs}.
  \item $d=2$, $r_1+r_2=5$, $|r_1-r_2|=1$. $|r_1-r_2|<d<r_1+r_2$~: \textbf{s\'ecants}.
  \item $\mathcal{C}_1$~: centre $O$, $r_1=4$~; $\mathcal{C}_2$~: centre $(3,0)$, $r_2=1$. $d=3$.
        $|r_1-r_2|=3=d$~: \textbf{tangents int\'erieurement}.
  \item $\mathcal{C}_1$~: centre $(1,1)$, $r_1=2\sqrt{2}$~; $\mathcal{C}_2$~: centre $(4,5)$, $r_2=\sqrt{2}$.
        $d=\sqrt{9+16}=5$~; $r_1+r_2=3\sqrt{2}\approx4{,}24$. $d>r_1+r_2$~: \textbf{ext\'erieurs}.
\end{enumerate}

%────────────────────────────────────────────────────────────────
\solutiontitle{6}{22}{Tangente ext\'erieure \`a un cercle depuis un point}

$\mathcal{C}$~: $x^2+y^2=9$~; $P(5,0)$.
\begin{enumerate}
  \item $5^2+0^2=25>9$~: $P$ est ext\'erieur.\checkmark
  \item Tangente $y=m(x-5)$, soit $mx-y-5m=0$.
        $d(O,\Delta)=\frac{|{-5m}|}{\sqrt{m^2+1}}=\frac{5|m|}{\sqrt{m^2+1}}$.
  \item $\frac{5|m|}{\sqrt{m^2+1}}=3\Rightarrow25m^2=9(m^2+1)\Rightarrow16m^2=9\Rightarrow m=\pm\frac{3}{4}$.
  \item Tangentes~: $y=\frac{3}{4}(x-5)$ et $y=-\frac{3}{4}(x-5)$.
\end{enumerate}

%────────────────────────────────────────────────────────────────
\solutiontitle{6}{23}{\'Equation cart\'esienne d'une parabole}

$F(0,1)$, directrice $y=-1$.
\begin{enumerate}
  \item $MF=\sqrt{x^2+(y-1)^2}$~; $d(M,y=-1)=y+1$.
  \item $x^2+(y-1)^2=(y+1)^2\Rightarrow x^2+y^2-2y+1=y^2+2y+1$.
  \item $x^2=4y$~: parabole de sommet $O$.
  \item Parabole ouverte vers le haut, sommet \`a l'origine.
\end{enumerate}

%────────────────────────────────────────────────────────────────
\solutiontitle{6}{24}{Droite d\'efinie par ses intercepts}

\begin{enumerate}
  \item $A(a,0)$~: $\frac{a}{a}+\frac{0}{b}=1$.\checkmark\quad $B(0,b)$~: $\frac{0}{a}+\frac{b}{b}=1$.\checkmark
  \item $a=3$, $b=4$~: $\frac{x}{3}+\frac{y}{4}=1\Rightarrow4x+3y-12=0$.
  \item $\frac{x}{4}+\frac{y}{-3}=1\Rightarrow-3x+4y+12=0$, soit $3x-4y=12$.
  \item $\frac{6}{a}+\frac{2}{4}=1\Rightarrow\frac{6}{a}=\frac{1}{2}\Rightarrow a=12$.
\end{enumerate}

%────────────────────────────────────────────────────────────────
\solutiontitle{6}{25}{Position relative~: syst\`eme}

\begin{enumerate}
  \item M\^eme pente $2$, $p_1\ne p_2$~: \textbf{parall\`eles}.
  \item $6x-2y+4=2(3x-y+2)$~: \textbf{confondues}.
  \item Pentes diff\'erentes ($1$ et $-1$)~: \textbf{s\'ecantes}.
  \item Pentes~: $d_1$ pente $-2/3$~; $d_2$ pente $-2/3$ mais $4x+6y=11\ne2(2x+3y-6)$. \textbf{Parall\`eles}.
\end{enumerate}

%────────────────────────────────────────────────────────────────
\solutiontitle{6}{26}{Cercle~: construire l'\'equation}

\begin{enumerate}
  \item $r=|7-3|=4$~: $(x-3)^2+(y+1)^2=16$.
  \item Milieu $\Omega=(1,3)$, $r=\frac{AB}{2}=\frac{\sqrt{16+4}}{2}=\frac{\sqrt{20}}{2}=\sqrt{5}$~:
        $(x-1)^2+(y-3)^2=5$.
  \item $r=\sqrt{25+144}=13$~: $x^2+y^2=169$.
  \item Centre $(2,3)$, tangent \`a $y=0$~: $r=|y_\Omega|=3$~: $(x-2)^2+(y-3)^2=9$.
\end{enumerate}

%────────────────────────────────────────────────────────────────
\solutiontitle{6}{27}{Calculer la distance~: applications}

$\mathcal{C}$~: $(x-2)^2+(y+1)^2=25$.
\begin{enumerate}
  \item $P_1(6,2)$~: $(6-2)^2+(2+1)^2=16+9=25$~: \textbf{sur} le cercle.
        $P_2(5,3)$~: $(5-2)^2+(3+1)^2=9+16=25$~: \textbf{sur} le cercle.
        $P_3(0,0)$~: $(0-2)^2+(0+1)^2=4+1=5<25$~: \textbf{int\'erieur}.
  \item $d(O,d)=\frac{|12|}{5}=\frac{12}{5}<3$~: la droite est \textbf{s\'ecante}.
  \item $3x+4y=12$, distance depuis $O$~: $12/5$. Corde de demi-longueur $\sqrt{9-144/25}=\sqrt{81/25}=9/5$.
        Points~: substituer $y=(12-3x)/4$ dans $x^2+y^2=9$. (Exercice de v\'erification num\'erique.)
\end{enumerate}

%────────────────────────────────────────────────────────────────
\solutiontitle{6}{28}{Triangle rectangle inscrit dans un cercle}

$\Omega(1,2)$, $r=4$.
\begin{enumerate}
  \item $A(-3,2)$~: $(-3-1)^2+(2-2)^2=16=4^2$.\checkmark\quad
        $B(5,2)$~: $(5-1)^2+(2-2)^2=16=4^2$.\checkmark\quad
        $C(1,6)$~: $(1-1)^2+(6-2)^2=16=4^2$.\checkmark
  \item Milieu de $[AB]$~: $\bigl(\frac{-3+5}{2},\frac{2+2}{2}\bigr)=(1,2)=\Omega$
        et $AB=8=2r$~: \textbf{oui, $[AB]$ est un diam\`etre}.
  \item $\vv{CA}=(-4,-4)$, $\vv{CB}=(4,-4)$.
        $\vv{CA}\cdot\vv{CB}=(-4)(4)+(-4)(-4)=-16+16=0$.\checkmark\quad $\widehat{ACB}=90^{\circ}$.\hfill$\square$
\end{enumerate}

%────────────────────────────────────────────────────────────────
\solutiontitle{6}{29}{R\'egion du plan d\'efinie par un cercle}

$\mathcal{C}$~: $x^2+y^2=9$.
\begin{enumerate}
  \item $\{M\mid OM<3\}$~: \textbf{int\'erieur ouvert} du disque.
  \item $\{M\mid OM\geq3\}$~: \textbf{ext\'erieur ferm\'e} (cercle inclus).
  \item $d(O, x+y=5)=\frac{5}{\sqrt{2}}\approx3{,}54>3$~: la droite est ext\'erieure au disque.
  \item $y=x+k$, forme $x-y+k=0$. $d=\frac{|k|}{\sqrt{2}}=3\Rightarrow k=\pm3\sqrt{2}$.
        Droites tangentes~: $y=x\pm3\sqrt{2}$.
\end{enumerate}

%────────────────────────────────────────────────────────────────
\solutiontitle{6}{30}{Trouver le cercle par conditions}

\begin{enumerate}
  \item Centre sur $Ox$~: $\Omega(a,0)$. $\Omega A^2=\Omega B^2$~:
        $a^2+9=(a-4)^2+1\Rightarrow a^2+9=a^2-8a+17\Rightarrow8a=8\Rightarrow a=1$.
        $r^2=1+9=10$~: cercle $(x-1)^2+y^2=10$.
  \item Centre sur $y=x$~: $\Omega(a,a)$. Tangent aux axes~: $r=|a|$.
        \'Equation~: $(x-a)^2+(y-a)^2=a^2$ (deux familles : $a>0$ et $a<0$).
  \item Centre sur $y=x$~: $\Omega(a,a)$. Passe par $A(2,0)$ et $B(-2,0)$~:
        $\Omega A=\Omega B\Rightarrow(a-2)^2+a^2=(a+2)^2+a^2\Rightarrow-8a=8a\Rightarrow a=0$.
        Centre $(0,0)$, $r=2$~: $x^2+y^2=4$ (unique).
\end{enumerate}

%────────────────────────────────────────────────────────────────
\solutiontitle{6}{31}{Droite et cercle~: intersection explicite}

\begin{enumerate}
  \item $3x+4y=0\Rightarrow y=-3x/4$. $x^2+9x^2/16=25\Rightarrow x^2=16\Rightarrow x=\pm4$.
        Points~: $(4,-3)$ et $(-4,3)$.
  \item $y=x-2$~: $(x-2)^2+(x-2)^2=16\Rightarrow2(x-2)^2=16\Rightarrow x-2=\pm2\sqrt{2}$.
        $x=2\pm2\sqrt{2}$~; points~: $(2+2\sqrt{2},2\sqrt{2})$ et $(2-2\sqrt{2},-2\sqrt{2})$.
  \item $y=1$~: $x^2+1-2x=0\Rightarrow(x-1)^2=0\Rightarrow x=1$.
        Point unique~: $(1,1)$ (tangente).
\end{enumerate}

%────────────────────────────────────────────────────────────────
\solutiontitle{6}{32}{Angle au centre vs angle inscrit}

Cercle unit\'e~; sommets \`a $\theta=\pi/6$, $5\pi/6$, $3\pi/2$.
\begin{enumerate}
  \item $A=(\cos(\pi/6),\sin(\pi/6))=(\frac{\sqrt{3}}{2},\frac{1}{2})$~;
        $B=(\cos(5\pi/6),\sin(5\pi/6))=(-\frac{\sqrt{3}}{2},\frac{1}{2})$~;
        $C=(\cos(3\pi/2),\sin(3\pi/2))=(0,-1)$.
  \item $AB=\sqrt{3}$~; $AC=\sqrt{3/4+9/4}=\sqrt{3}$~; $BC=\sqrt{3}$.\checkmark\quad \'Equilat\'eral.
  \item P\'erim\`etre~: $3\sqrt{3}$~; Aire~: $\frac{\sqrt{3}}{4}(\sqrt{3})^2=\frac{3\sqrt{3}}{4}$.
\end{enumerate}

%────────────────────────────────────────────────────────────────
\solutiontitle{6}{33}{Lieu g\'eom\'etrique~: exercice suppl\'ementaire}

$MA/MB=2$, $A(0,0)$, $B(3,0)$.
\begin{enumerate}
  \item $x^2+y^2=4[(x-3)^2+y^2]\Rightarrow x^2+y^2=4x^2-24x+36+4y^2$.
        $3x^2+3y^2-24x+36=0\Rightarrow x^2+y^2-8x+12=0$.
  \item $(x-4)^2+y^2=16-12=4$.
  \item Centre $(4,0)$, rayon $2$ (cercle d'Apollonius).
\end{enumerate}

%────────────────────────────────────────────────────────────────
\solutiontitle{6}{34}{\logicoalgo\ Algorithme~: position d'une droite et d'un cercle}

\begin{enumerate}
  \item $d>r$, $d=r$, $d<r$ forment une \textbf{partition} de $\R^+$~:
        $d\in\R^+$ et exactement l'une des trois conditions tient.
        Logiquement~: $(d>r)\vee(d=r)\vee(d<r)$, termes deux \`a deux exclusifs.
  \item $\Omega(1,-2)$, $r=5$, droite $3x-4y+2=0$~:
        $d=\frac{|3-4(-2)+2|}{5}=\frac{13}{5}=2{,}6<5$~: \textbf{s\'ecante}.
  \item $\Omega(0,0)$, $r=3$, droite $y=x+4$ (soit $x-y+4=0$)~:
        $d=\frac{4}{\sqrt{2}}=2\sqrt{2}\approx2{,}83<3$~: \textbf{s\'ecante}.
  \item Si $d=r$~: le point de contact $H$ est le projet\'e de $\Omega$ sur $\Delta$.
        $H=\Omega-d\cdot\hat{n}$ o\`u $\hat{n}$ est le vecteur normal unitaire.
\end{enumerate}

%────────────────────────────────────────────────────────────────
\solutiontitle{6}{35}{Cercle de diam\`etre $[AB]$}

$A(a,0)$, $B(-a,0)$, $M(x,y)$.
\begin{enumerate}
  \item $\vv{MA}=(a-x,-y)$~; $\vv{MB}=(-a-x,-y)$.
        $\vv{MA}\cdot\vv{MB}=(a-x)(-a-x)+y^2=-(a-x)(a+x)+y^2=-(a^2-x^2)+y^2=x^2+y^2-a^2$.
  \item $\vv{MA}\cdot\vv{MB}=0\Leftrightarrow x^2+y^2=a^2$.
  \item C'est le cercle de centre $O$ et rayon $a=\dfrac{AB}{2}$~: cercle de diam\`etre $[AB]$.\hfill$\square$
\end{enumerate}

%────────────────────────────────────────────────────────────────
\solutiontitle{6}{36}{Distance entre droites parall\`eles}

\begin{enumerate}
  \item Prendre $P(x_0,y_0)$ sur la premi\`ere droite ($ax_0+by_0+c_1=0$).
        $d(P,\Delta_2)=\frac{|ax_0+by_0+c_2|}{\sqrt{a^2+b^2}}=\frac{|c_2-c_1|}{\sqrt{a^2+b^2}}$.\hfill$\square$
  \item $d=\frac{|5-(-10)|}{5}=\frac{15}{5}=3$.
  \item $y=2x+1$ et $y=2x-4$~: forme $2x-y+1=0$ et $2x-y-4=0$.
        $d=\frac{|1-(-4)|}{\sqrt{5}}=\frac{5}{\sqrt{5}}=\sqrt{5}$.
\end{enumerate}

%────────────────────────────────────────────────────────────────
\solutiontitle{6}{37}{Bipoint et axe radical}

\begin{enumerate}
  \item $\Omega_1A=\Omega_1B=r_1$ et $\Omega_2A=\Omega_2B=r_2$~:
        les deux centres sont \'equidistants de $A$ et $B$~; ils sont donc tous les deux sur la
        m\'ediatrice de $[AB]$. La droite $\Omega_1\Omega_2$ est la m\'ediatrice de $[AB]$,
        donc perpendiculaire \`a $AB$.\hfill$\square$
  \item Le milieu de $[AB]$ est sur la m\'ediatrice de $[AB]$, donc sur $\Omega_1\Omega_2$.\hfill$\square$
  \item $\mathcal{C}_1$~: $x^2+y^2=10$~; $\mathcal{C}_2$~: $(x-2)^2+y^2=6$, centre $(2,0)$.
        Soustraction~: $x^2-(x-2)^2=4\Rightarrow4x-4=4\Rightarrow x=2$.
        En $x=2$~: $y^2=10-4=6$~; $A=(2,\sqrt{6})$, $B=(2,-\sqrt{6})$.
        Corde $AB$~: $x=2$ (verticale). Droite des centres~: $y=0$ (horizontale).\checkmark\quad Perp.
        Milieu de $[AB]$~: $(2,0)$~: sur $y=0$.\checkmark
\end{enumerate}

%────────────────────────────────────────────────────────────────
\solutiontitle{6}{38}{\logiq\ Propositions sur les droites}

\begin{enumerate}
  \item \textbf{Fausse} si on inclut les droites verticales (pas de pente d\'efinie).
        Si $m_1=m_2$ et les droites sont non verticales~: vraie.
        N\'egation~: \og Il existe deux droites de m\^eme pente qui ne sont pas parall\`eles.\fg
  \item \textbf{Vraie.} N\'egation~: \og $\exists$ deux points par lesquels aucune droite ne passe.\fg
  \item \textbf{Vraie.} N\'egation~: \og $\exists$ une droite et un point tels que la perpendiculaire
        par ce point n'est pas unique.\fg
  \item \textbf{Fausse.} Deux droites perpendiculaires se coupent (en $90^{\circ}$).
        N\'egation~: \og Il existe des droites perpendiculaires qui se coupent.\fg (trivial).
\end{enumerate}

%────────────────────────────────────────────────────────────────
\solutiontitle{6}{39}{\logiq\ \'Equivalences sur les droites}

\begin{enumerate}
  \item $\Leftrightarrow$~: para. ssi m\^eme pente (pour les non-verticales).
  \item $\Leftrightarrow$~: d\'efinition de la perpendicularit\'e via les pentes.
  \item $\Leftrightarrow$~: un point est sur la droite ssi ses coordonn\'ees satisfont l'\'equation.
  \item $\Leftrightarrow$~: d\'efinition de la tangence.
\end{enumerate}

%────────────────────────────────────────────────────────────────
\solutiontitle{6}{40}{\logiq\ Raisonnement sur la position droite/cercle}

\begin{enumerate}
  \item Les trois cas forment une partition~: $d(\Omega,\Delta)\in\R^+$ v\'erifie
        exactement l'une de~: $>r$, $=r$, $<r$. Logiquement~:
        $\forall d$~: $(d>r)\veebar(d=r)\veebar(d<r)$.\hfill$\square$
  \item Supposons deux points de contact $T_1$ et $T_2$. Alors $d(\Omega,\Delta)\leq\Omega T_i=r$
        pour $i=1,2$, mais la distance du centre \`a la droite est unique. Contradiction.\hfill$\square$
  \item Contrapos\'ee~: si $d(\Omega,\Delta)\ne r$, la droite n'est pas tangente
        (elle est soit s\'ecante, soit ext\'erieure).
\end{enumerate}

%────────────────────────────────────────────────────────────────
\solutiontitle{6}{41}{\logiq\ Lieux g\'eom\'etriques et \'equivalences}

\begin{enumerate}
  \item $M\in$ m\'ediatrice de $[AB]\Leftrightarrow MA=MB$.
  \item $M\in$ cercle de diam\`etre $[AB]\Leftrightarrow\vv{MA}\cdot\vv{MB}=0$.
  \item $M\in\mathcal{C}(\Omega,r)\Leftrightarrow\Omega M=r$.
  \item (3)~: $\Omega M=r\Leftrightarrow\Omega M^2=r^2\Leftrightarrow(x-a)^2+(y-b)^2=r^2$.\hfill$\square$
\end{enumerate}

%────────────────────────────────────────────────────────────────
\solutiontitle{6}{42}{\logiq\ Conditions sur l'\'equation d'un cercle}

\begin{enumerate}
  \item La condition $D^2+E^2-4F>0$ est une \textbf{CNS} pour que l'\'equation repr\'esente un cercle.
  \item N\'egation~: \og $D^2+E^2-4F\leq0$\fg\ (point d\'eg\'en\'er\'e ou ensemble vide).
  \item Exemple~: $x^2+y^2=0$ (point unique $O$) : $D=E=F=0$~; $D^2+E^2-4F=0$.
  \item Exemple~: $x^2+y^2+1=0$~: $F=1$, $D=E=0$~; $-4<0$~: ensemble vide.
\end{enumerate}

%────────────────────────────────────────────────────────────────
\solutiontitle{6}{43}{\logiq\ Th\'eor\`eme de Thal\`es~: logique}

\begin{enumerate}
  \item $P$~: \og $M\in[AB]$ et $AM/AB=k$ et droite par $M\parallel BC$\fg.
        $Q$~: \og La parallèle coupe $[AC]$ en $N$ avec $AN/AC=k$.\fg
        Th\'eor\`eme~: $P\Rightarrow Q$.
  \item R\'eciproque vraie (th\'eor\`eme r\'eciproque de Thal\`es)~:
        si $AN/AC=AM/AB$, la droite $MN$ est parall\`ele \`a $BC$.
  \item Contrapos\'ee~: si $AN/AC\ne AM/AB$, alors $MN\not\parallel BC$.
  \item Avec $A(0,0)$, $B(6,0)$, $C(0,4)$~: voir Prob. 11 du Ch5 (m\^eme configuration).
\end{enumerate}

%────────────────────────────────────────────────────────────────
\solutiontitle{6}{44}{\algo\ Algorithme~: \'equation d'une droite par deux points}

Traces~:
$(0,0),(3,6)$~: $m=2$, $p=0$~; $y=2x$.
$(2,1),(2,5)$~: $x_A=x_B=2$~; droite verticale $x=2$.
$(-1,4),(3,-2)$~: $m=\frac{-6}{4}=-\frac{3}{2}$, $p=4-(-\frac{3}{2})(-1)=4-\frac{3}{2}=\frac{5}{2}$~; $y=-\frac{3}{2}x+\frac{5}{2}$.

\begin{lstlisting}[language=Python]
def droite(xA, yA, xB, yB):
    if xA == xB:
        print(f"x = {xA}")
    else:
        m = (yB - yA) / (xB - xA)
        p = yA - m * xA
        print(f"y = {m}x + {p}")
\end{lstlisting}

%────────────────────────────────────────────────────────────────
\solutiontitle{6}{45}{\algo\ Algorithme~: position relative droite/cercle \computer}

\begin{lstlisting}[language=Python]
import math

def position_et_intersections(a, b, r, m, p):
    # droite: y = mx + p, soit mx - y + p = 0
    dist = abs(m*a - b + p) / math.sqrt(m**2 + 1)
    if dist > r:
        print("Exterieure")
    elif abs(dist - r) < 1e-9:
        print("Tangente")
    else:
        print("Secante")
        # substitution y = mx + p dans (x-a)^2 + (y-b)^2 = r^2
        A = 1 + m**2
        B = 2*(m*(p-b) - a)
        C = a**2 + (p-b)**2 - r**2
        delta = B**2 - 4*A*C
        x1 = (-B - math.sqrt(delta)) / (2*A)
        x2 = (-B + math.sqrt(delta)) / (2*A)
        print(f"  ({x1:.4f}, {m*x1+p:.4f}) et ({x2:.4f}, {m*x2+p:.4f})")

# Test: C: x^2+y^2=25 (centre (0,0), r=5), d: y=x+1
position_et_intersections(0, 0, 5, 1, 1)
\end{lstlisting}

%────────────────────────────────────────────────────────────────
\solutiontitle{6}{46}{\algo\ Algorithme~: m\'ediatrice \computer}

\begin{lstlisting}[language=Python]
def mediatrice(xA, yA, xB, yB):
    Ix, Iy = (xA+xB)/2, (yA+yB)/2
    dx, dy = xB-xA, yB-yA
    # equation: dx*(x - Ix) + dy*(y - Iy) = 0
    # soit: dx*x + dy*y = dx*Ix + dy*Iy
    c = dx*Ix + dy*Iy
    print(f"Mediatrice: {dx}x + {dy}y = {c}")
    # Verification: le milieu est sur la mediatrice
    verif = dx*Ix + dy*Iy
    print(f"Verif milieu: {verif} = {c}", verif == c)

mediatrice(1, 3, 5, 1)
\end{lstlisting}

%────────────────────────────────────────────────────────────────
\solutiontitle{6}{47}{\algo\ Algorithme~: tangente \`a un cercle depuis un point ext\'erieur \computer}

\begin{lstlisting}[language=Python]
import math

def tangentes_exterieures(a, b, r, px, py):
    # Tangentes y - py = m(x - px), soit m*x - y + (py - m*px) = 0
    # d(centre, droite) = |m*a - b + py - m*px| / sqrt(m^2+1) = r
    # => (m*(a-px) + (py-b))^2 = r^2*(m^2+1)
    A = (a-px)**2 - r**2
    B = 2*(a-px)*(py-b)
    C = (py-b)**2 - r**2
    delta = B**2 - 4*A*C
    if delta < 0:
        print("Pas de tangente reelle")
        return
    m1 = (-B + math.sqrt(delta)) / (2*A)
    m2 = (-B - math.sqrt(delta)) / (2*A)
    for m in [m1, m2]:
        p = py - m*px
        print(f"y = {m:.4f}x + {p:.4f}")

tangentes_exterieures(0, 0, 3, 5, 0)
# Attendu: m = +/-3/4
\end{lstlisting}

%────────────────────────────────────────────────────────────────
\solutiontitle{6}{48}{\algo\ Algorithme~: cercle passant par trois points \computer}

\begin{lstlisting}[language=Python]
import numpy as np

def cercle_3pts(A, B, C):
    # x^2+y^2+Dx+Ey+F = 0
    M = np.array([[A[0],A[1],1],[B[0],B[1],1],[C[0],C[1],1]], float)
    r = np.array([-(A[0]**2+A[1]**2),
                  -(B[0]**2+B[1]**2),
                  -(C[0]**2+C[1]**2)], float)
    D, E, F = np.linalg.solve(M, r)
    cx, cy = -D/2, -E/2
    rad = math.sqrt(cx**2 + cy**2 - F)
    print(f"Centre ({cx:.4f}, {cy:.4f}), rayon {rad:.4f}")

import math
cercle_3pts((0,0),(4,0),(1,3))
\end{lstlisting}

%────────────────────────────────────────────────────────────────
\solutiontitle{6}{49}{\algo\ Algorithme~: distance point--droite sur une grille \computer}

\begin{lstlisting}[language=Python]
import math

# Droite y = 2x+1, soit 2x - y + 1 = 0
a, b, c = 2, -1, 1
for x in range(-5, 6):
    for y in range(-5, 6):
        d = abs(a*x + b*y + c) / math.sqrt(a**2 + b**2)
        if d < 1:
            print(f"({x},{y}): d = {d:.4f}")
\end{lstlisting}

%────────────────────────────────────────────────────────────────
\subsection*{Probl\`emes}
\addcontentsline{toc}{subsection}{Probl\`emes}

%────────────────────────────────────────────────────────────────
\psolutiontitle{6}{1}{Triangle inscrit dans un cercle}

$A(5,0)$, $B(-5,0)$, $C(0,5)$.
\textbf{(1)} $25+0=25$.\checkmark\quad $25+0=25$.\checkmark\quad $0+25=25$.\checkmark
\textbf{(2)} $AB=10$~; $AC=\sqrt{25+25}=5\sqrt{2}$~; $BC=5\sqrt{2}$.
\textbf{(3)} $\vv{CA}=(5,-5)$, $\vv{CB}=(-5,-5)$.
$\cos\hat{C}=\frac{\vv{CA}\cdot\vv{CB}}{|CA||CB|}=\frac{-25+25}{50}=0$~; $\hat{C}=90^{\circ}$.
$\hat{A}=\hat{B}=45^{\circ}$ (triangle isoc\`ele rectangle).
\textbf{(4)} Aire~: $\frac{1}{2}|det(\vv{AB},\vv{AC})|=\frac{1}{2}|(-10)(5)-(0)(-5)|=25$.
Base $\times$ hauteur~: $\frac{1}{2}\cdot10\cdot5=25$.\checkmark

%────────────────────────────────────────────────────────────────
\psolutiontitle{6}{2}{Intersection cercle--droite et corde}

$\mathcal{C}$~: $x^2+y^2=25$~; $(\Delta)$~: $3x+4y-25=0$.
\textbf{(1)} $d(O,\Delta)=\frac{25}{5}=5=r$~: \textbf{tangente}.
\textbf{(2)} Point de tangence~: projet\'e de $O$ sur $\Delta$.
$H=O+t(3,4)$~: $3(3t)+4(4t)-25=0\Rightarrow25t=25\Rightarrow t=1$.
$H=(3,4)$.
\textbf{(3)} Corde~: tangente, donc 0 corde (longueur 0).
\textbf{(4)} Tangente en $(3,4)$~: $3x+4y=25$.\checkmark (c'est la droite elle-m\^eme).

%────────────────────────────────────────────────────────────────
\psolutiontitle{6}{3}{M\'ediane, hauteur, m\'ediatrice~: r\'ecap}

$A(2,6)$, $B(-4,0)$, $C(4,0)$.
\textbf{(1)} centre de gravité~: $G=(\frac{2}{3},2)$.
Circocentre~: m\'ediatrice de $BC$~: $x=0$~; m\'ediatrice de $AB$~: milieu $(-1,3)$, pente $AB=1$, perp. $-1$~: $y=-x+2$. En $x=0$~: $y=2$. Mais $O'A=\sqrt{4+16}=\sqrt{20}$, $O'B=\sqrt{4+4}=\sqrt{8}$~: $O'A\ne O'B$.
\emph{Correction~: milieu de $AB=(-1,3)$, pente de $AB=\frac{6}{6}=1$, perp. $m=-1$~: $y-3=-(x+1)\Rightarrow y=-x+2$. En $x=0$~: $y=2$. $O'=(0,2)$.
$O'A=\sqrt{4+16}=\sqrt{20}$~; $O'C=\sqrt{16+4}=\sqrt{20}$~; $O'B=\sqrt{16+4}=\sqrt{20}$.\checkmark}
Orthocentre~: hauteur de $A\perp BC$~: $x=2$.
Hauteur de $B\perp AC$~: pente $AC=\frac{-6}{2}=-3$, perp. $\frac{1}{3}$, par $B(-4,0)$~: $y=\frac{x+4}{3}$.
En $x=2$~: $y=2$. $H=(2,2)$.

\textbf{(2)} $O'=(0,2)$, $G=(\frac{2}{3},2)$, $H=(2,2)$~: m\^eme ordonn\'ee $\Rightarrow$ align\'es sur $y=2$.\checkmark
$\vv{O'G}=(\frac{2}{3},0)$~; $\vv{GH}=(\frac{4}{3},0)=2\vv{O'G}$~:
$O'G=\dfrac{GH}{2}$.\checkmark

\textbf{(3)} $AB=AC=\sqrt{36+36}=6\sqrt{2}$~: isoc\`ele. $BC=8$.
$\vv{AB}=(-6,-6)$, $\vv{AC}=(2,-6)$~: $\cos\hat{A}=\frac{-12+36}{6\sqrt{2}\cdot2\sqrt{10}}<0$~: obtusangle.

%────────────────────────────────────────────────────────────────
\psolutiontitle{6}{4}{Le th\'eor\`eme de Thal\`es en coordonn\'ees}

$A(0,0)$, $B(6,0)$, $C(0,4)$. $M$ sur $[AB]$ avec $AM=4$~: $M=(4,0)$.
\textbf{(1)} $M=(4,0)$ (point sur l'axe $x$).
\textbf{(2)} Droite parall\`ele \`a $BC$ passant par $M$~: $BC$ pente $\frac{-4}{6}=-\frac{2}{3}$.
Droite~: $y=-\frac{2}{3}(x-4)=-\frac{2}{3}x+\frac{8}{3}$.
$[AC]$~: $x=0$~; $N=(0,\frac{8}{3})$.
\textbf{(3)} $AM/AB=4/6=2/3$~; $AN/AC=(8/3)/4=2/3$.\checkmark
\textbf{(4)} $MN=\sqrt{16+64/9}=\sqrt{208/9}=\frac{4\sqrt{13}}{3}$~; $BC=\sqrt{36+16}=\sqrt{52}=2\sqrt{13}$.
$MN/BC=\frac{4\sqrt{13}/3}{2\sqrt{13}}=\frac{2}{3}=AM/AB$.\checkmark

%────────────────────────────────────────────────────────────────
\psolutiontitle{6}{5}{Triangle quelconque inscrit}

$A(2,0)$, $B(-2,0)$, $C(0,3)$.
\textbf{(1)} Syst\`eme $3\times3$~: $F=0$ (par $A$ et $B$ sym\'etriques), $D=0$ (sym\'etrie).
$4+2E+F=0$~: non, $B(-2,0)$~: $4-2D+F=0$. $4+2D+F=0$. D'o\`u $D=0$, $F=-4$.
Par $C(0,3)$~: $9+3E-4=0\Rightarrow E=-5/3$.
Cercle~: $x^2+y^2-\frac{5}{3}y-4=0\Rightarrow x^2+(y-\frac{5}{6})^2=4+\frac{25}{36}=\frac{169}{36}$.
Centre $\Omega=(0,\frac{5}{6})$, rayon $r=\frac{13}{6}$.
\textbf{(2)} $AB=4$~; $AC=\sqrt{4+9}=\sqrt{13}$~; $BC=\sqrt{4+9}=\sqrt{13}$. P\'erim\`etre~: $4+2\sqrt{13}$.
\textbf{(3)} Par le produit scalaire~: $\hat{A}$~: $\vv{AB}=(-4,0)$, $\vv{AC}=(-2,3)$.
$\cos\hat{A}=\frac{8}{4\sqrt{13}}=\frac{2}{\sqrt{13}}\Rightarrow\hat{A}\approx56{,}3^{\circ}$.
$\hat{B}=\hat{A}\approx56{,}3^{\circ}$ (isoc\`ele). $\hat{C}\approx67{,}4^{\circ}$.

%────────────────────────────────────────────────────────────────
\psolutiontitle{6}{6}{Droites et syst\`emes~: triangle}

$d_1$~: $y=2x$~; $d_2$~: $y=-x+6$~; $d_3$~: $y=1$.
\textbf{(1)} $d_1\cap d_2$~: $2x=-x+6\Rightarrow x=2$~; $A=(2,4)$.
$d_1\cap d_3$~: $y=1\Rightarrow x=1/2$~; $B=(1/2,1)$.
$d_2\cap d_3$~: $1=-x+6\Rightarrow x=5$~; $C=(5,1)$.
\textbf{(2)} $AB=\sqrt{(3/2)^2+9}=\sqrt{9/4+9}=\frac{3\sqrt{5}}{2}$~;
$AC=\sqrt{9+9}=3\sqrt{2}$~; $BC=9/2$.
Aire~: $\frac{1}{2}\cdot BC\cdot h=\frac{1}{2}\cdot\frac{9}{2}\cdot3=\frac{27}{4}$.
\textbf{(3)} Le cercle inscrit a son centre \`a \'egale distance des trois c\^ot\'es.
Distance de $I(x_I,y_I)$ \`a $d_3$~: $|y_I-1|$~;
\`a $d_1$~: $\frac{|2x_I-y_I|}{\sqrt{5}}$~; \`a $d_2$~: $\frac{|x_I+y_I-6|}{\sqrt{2}}$.
(Le calcul complet du cercle inscrit g\'en\'eral d\'epasse le niveau Seconde.)

%────────────────────────────────────────────────────────────────
\psolutiontitle{6}{7}{Perpendiculaire en un point}

$(\Delta)$~: $2x-3y+6=0$~; $A(4,5)$.
\textbf{(1)} La pente de $(\Delta)$ est $2/3$ ; la perpendiculaire a donc
pour pente $-3/2$. Par $A$~:
$y-5=-\frac32(x-4)$, soit $3x+2y-22=0$.
\textbf{(2)} En r\'esolvant $2x-3y+6=0$ et $3x+2y-22=0$, on obtient
$H=\left(\frac{54}{13},\frac{62}{13}\right)$.
\textbf{(3)} $AH=\sqrt{\left(\frac{2}{13}\right)^2+\left(\frac{3}{13}\right)^2}
=\frac{\sqrt{13}}{13}=\frac{1}{\sqrt{13}}$.
\textbf{(4)} $d(A,\Delta)=\frac{|2(4)-3(5)+6|}{\sqrt{2^2+(-3)^2}}
=\frac{1}{\sqrt{13}}=AH$.\checkmark

%────────────────────────────────────────────────────────────────
\psolutiontitle{6}{8}{Syst\`eme de cercles co\-axiaux}

$\mathcal{C}_1$~: $x^2+y^2-4x=0\Rightarrow(x-2)^2+y^2=4$. Centre $(2,0)$, rayon $2$.
$\mathcal{C}_2$~: $x^2+y^2+2x-6=0\Rightarrow(x+1)^2+y^2=7$. Centre $(-1,0)$, rayon $\sqrt{7}$.
\textbf{(2)} $d=3$~; $r_1+r_2=2+\sqrt{7}\approx4{,}65$~; $|r_1-r_2|=\sqrt{7}-2\approx0{,}65$.
$|r_1-r_2|<d<r_1+r_2$~: \textbf{s\'ecants}.
\textbf{(3)} Soustraction~: $(x^2+y^2-4x)-(x^2+y^2+2x-6)=0\Rightarrow-6x+6=0\Rightarrow x=1$.
En $x=1$~: $1+y^2-4=0\Rightarrow y^2=3$~; $A=(1,\sqrt{3})$, $B=(1,-\sqrt{3})$.
\textbf{(4)} Axe radical~: $-6x+6=0\Rightarrow x=1$.\checkmark

%────────────────────────────────────────────────────────────────
\psolutiontitle{6}{9}{Localisation par distance~: GPS simplifi\'e}

$A(0,0)$, $B(10,0)$, $C(5,8)$~; $d_A=d_B=5$, $d_C=8$.
\textbf{(1)} $x^2+y^2=25$.
\textbf{(2)} $(x-10)^2+y^2=25$~; $(x-5)^2+(y-8)^2=64$.
\textbf{(3)} Cercles $A$ et $B$~: $x^2+y^2=(x-10)^2+y^2\Rightarrow20x=100\Rightarrow x=5$.
En $x=5$~: $25+y^2=25\Rightarrow y=0$. $M=(5,0)$.
V\'erif. avec $C$~: $(5-5)^2+(0-8)^2=64=8^2$.\checkmark
\textbf{(4)} Les cercles $A$ et $B$ donnent deux solutions~: $(5,0)$ et $(5,0)$... en fait une seule
ici car $A$ et $B$ sont sym\'etriques. En g\'en\'eral deux cercles donnent $0$, $1$ ou $2$ points.
Le troisi\`eme cercle (centre $C$, rayon $d_C$) ne passe que par l'un des candidats~:
il permet de s\'electionner la solution unique.\hfill$\square$

%────────────────────────────────────────────────────────────────
\psolutiontitle{6}{10}{Triangle rectangle et cercle}

\textbf{(1)} Si $\vv{CA}\cdot\vv{CB}=0$, alors par le th\'eor\`eme du cercle de diam\`etre
(prouv\'e en exo 35)~: $C$ est sur le cercle de diam\`etre $[AB]$.\hfill$\square$
\textbf{(2)} R\'eciproque vraie~: si $C$ est sur le cercle de diam\`etre $[AB]$,
alors $\vv{CA}\cdot\vv{CB}=0$, donc $\hat{ACB}=90^{\circ}$.\hfill$\square$
\textbf{(3)} $A(-3,0)$, $B(3,0)$~; cercle~: $x^2+y^2=9$.
Isoc\`ele $AC=BC$~: $C$ sur l'axe $x=0$. Sur le cercle~: $y^2=9\Rightarrow C=(0,\pm3)$.
V\'erif.~: $\vv{CA}=(-3,-3)$, $\vv{CB}=(3,-3)$~; $-9+9=0$.\checkmark

%────────────────────────────────────────────────────────────────
\psolutiontitle{6}{11}{Distance minimale}

$(\Delta)$~: $y=2x+1$ ($2x-y+1=0$)~; $A(5,3)$.
\textbf{(1)} Perp. de pente $-1/2$ par $A$~: $y-3=-\frac{1}{2}(x-5)\Rightarrow x+2y-11=0$.
\textbf{(2)} Intersection~: $2x-y=-1$ et $x+2y=11$.
$4x-2y=-2$ et $x+2y=11$~: $5x=9\Rightarrow x=9/5$~; $y=\frac{11-9/5}{2}=\frac{46/5}{2}=\frac{23}{5}$.
$H=(9/5,23/5)$.
\textbf{(3)} $AH=\sqrt{(5-9/5)^2+(3-23/5)^2}=\sqrt{(16/5)^2+(8/5)^2}=\frac{\sqrt{320}}{5}=\frac{8\sqrt{5}}{5}$.
\textbf{(4)} $d(A,\Delta)=\frac{|10-3+1|}{\sqrt{5}}=\frac{8}{\sqrt{5}}=\frac{8\sqrt{5}}{5}$.\checkmark

%────────────────────────────────────────────────────────────────
\psolutiontitle{6}{12}{Angle au centre et arc}

$\mathcal{C}$~: centre $O$, rayon $5$~; $A(5,0)$, $B(0,5)$, $C(-5,0)$, $D(0,-5)$.
\textbf{(1)} $\vv{OA}\cdot\vv{OB}=0$~: $\widehat{AOB}=90^{\circ}$.
\textbf{(2)} $\widehat{AOC}=180^{\circ}$~; $\widehat{BOD}=180^{\circ}$.
\textbf{(3)} $M(3,4)$~: $9+16=25$.\checkmark\quad $\widehat{AOM}$~: $\cos\theta=\frac{15}{25}=\frac{3}{5}$~; $\theta=\arccos(3/5)\approx53{,}1^{\circ}$.
\textbf{(4)} P\'erim\`etre du carr\'e $ABCD$~: $4\cdot5\sqrt{2}=20\sqrt{2}\approx28{,}3$.
Circonf\'erence~: $2\pi\cdot5\approx31{,}4$. $20\sqrt{2}<10\pi$~: le carr\'e est inscrit \`a l'int\'erieur.

%────────────────────────────────────────────────────────────────
\psolutiontitle{6}{13}{\logiq\ Logique des cercles}

\textbf{(1)} $P$ int\'erieur \`a $\mathcal{C}(\Omega,r)\Leftrightarrow\Omega P<r\Leftrightarrow(x_P-a)^2+(y_P-b)^2<r^2$.
\textbf{(2)} $d(\Omega_1,\Omega_2)=r_1+r_2$~: le point de contact $T$ est sur $[\Omega_1\Omega_2]$
avec $\Omega_1T=r_1$ et $\Omega_2T=r_2$. $T$ est sur $\mathcal{C}_1$ et $\mathcal{C}_2$.\hfill$\square$
\textbf{(3)} N\'egation de \og s\'ecants\fg~: $|r_1-r_2|\geq d$ ou $d\geq r_1+r_2$
(ext\'erieurs ou l'un dans l'autre, y compris les cas tangents).

%────────────────────────────────────────────────────────────────
\psolutiontitle{6}{14}{\logiq\ Logique et lieux g\'eom\'etriques}

\textbf{(1)} Voir exo 35 (d\'emonstration compl\`ete).\hfill$\square$
\textbf{(2)} L'axe radical~: $\pi_1(M)=\pi_2(M)\Leftrightarrow(D_1-D_2)x+(E_1-E_2)y+(F_1-F_2)=0$.
C'est une \'equation lin\'eaire (les $D_i,E_i,F_i$ ne sont pas tous \'egaux si les cercles ne sont pas concentriques)~: c'est bien une droite.\hfill$\square$
\textbf{(3)} $\forall M\in\mathcal{A}$~: $M\Omega_1^2-r_1^2=M\Omega_2^2-r_2^2$.

%────────────────────────────────────────────────────────────────
\psolutiontitle{6}{15}{\algo\ GPS simplifi\'e \computer}

\begin{lstlisting}[language=Python]
import numpy as np

def gps(A, rA, B, rB, C, rC):
    # Systeme: (x-Ax)^2+(y-Ay)^2=rA^2, etc.
    # Soustraire A-B et A-C pour obtenir 2 equations lineaires
    M = np.array([
        [2*(B[0]-A[0]), 2*(B[1]-A[1])],
        [2*(C[0]-A[0]), 2*(C[1]-A[1])]
    ], float)
    r = np.array([
        rA**2 - rB**2 - A[0]**2 + B[0]**2 - A[1]**2 + B[1]**2,
        rA**2 - rC**2 - A[0]**2 + C[0]**2 - A[1]**2 + C[1]**2
    ], float)
    return np.linalg.solve(M, r)
\end{lstlisting}

%────────────────────────────────────────────────────────────────
\psolutiontitle{6}{16}{\algo\ Tracer une droite et un cercle \computer}

\begin{lstlisting}[language=Python]
import numpy as np, matplotlib.pyplot as plt, math

def tracer(a, b, r, m, p):
    theta = np.linspace(0, 2*np.pi, 300)
    plt.figure()
    plt.plot(a + r*np.cos(theta), b + r*np.sin(theta), 'b-', label='Cercle')
    x = np.linspace(a-r-1, a+r+1, 200)
    plt.plot(x, m*x + p, 'r-', label='Droite')
    d = abs(m*a - b + p) / math.sqrt(m**2 + 1)
    if d > r: pos = "Exterieure"
    elif abs(d-r) < 1e-9: pos = "Tangente"
    else:
        pos = "Secante"
        A_c = 1+m**2; B_c = 2*(m*(p-b)-a); C_c = a**2+(p-b)**2-r**2
        delta = B_c**2 - 4*A_c*C_c
        for xi in [(-B_c+math.sqrt(delta))/(2*A_c), (-B_c-math.sqrt(delta))/(2*A_c)]:
            plt.plot(xi, m*xi+p, 'ko', markersize=8)
    plt.title(pos); plt.legend(); plt.axis('equal'); plt.grid(True); plt.show()
\end{lstlisting}

%────────────────────────────────────────────────────────────────
\subsection*{D\'efis}
\addcontentsline{toc}{subsection}{D\'efis}

%────────────────────────────────────────────────────────────────
\noindent\phantomsection
{\color{BO}\bfseries\small D\'efi~1~\textendash\ \textit{Cercle circonscrit~: construction}}
\par\smallskip

$A(0,0)$, $B(4,0)$, $C(1,3)$.
Syst\`eme~: $F=0$~; $16+4D+F=0\Rightarrow D=-4$~; $1+9+D+3E+F=0\Rightarrow E=-2$.
\'Equation~: $(x-2)^2+(y-1)^2=5$.
Centre $(2,1)$, rayon $\sqrt{5}$.
V\'erif.~: $A$~: $4+1=5$.\checkmark\quad $B$~: $4+1=5$.\checkmark\quad $C$~: $1+4=5$.\checkmark

%────────────────────────────────────────────────────────────────
\noindent\phantomsection
{\color{BO}\bfseries\small D\'efi~2~\textendash\ \textit{Th\'eor\`eme de Ptol\'em\'ee}}
\par\smallskip

$A(1,0)$, $B(0,1)$, $C(-1,0)$, $D(0,-1)$~; cercle unit\'e.
$AC=2$, $BD=2$, $AB=BC=CD=DA=\sqrt{2}$.
$AC\cdot BD=4$~; $AB\cdot CD+AD\cdot BC=\sqrt{2}\cdot\sqrt{2}+\sqrt{2}\cdot\sqrt{2}=2+2=4$.\checkmark
Cercle~: $x^2+y^2=1$.

%────────────────────────────────────────────────────────────────
\noindent\phantomsection
{\color{BO}\bfseries\small D\'efi~3~\textendash\ \textit{Positions relatives~: discussion selon un param\`etre}}
\par\smallskip

$(\Delta_k)$~: $x+y=k$, soit $x+y-k=0$~; $\mathcal{C}$~: $x^2+y^2=r^2$.
$d=\frac{|k|}{\sqrt{2}}$.

S\'ecante~: $|k|<r\sqrt{2}$~; tangente~: $|k|=r\sqrt{2}$~; ext\'erieure~: $|k|>r\sqrt{2}$.

Pour $r=3$~: tangence si $|k|=3\sqrt{2}$, soit $k=\pm3\sqrt{2}$.
Points de tangence~: projet\'e de $O$ sur $\Delta_k$~:
$H=t(1,1)$ avec $t+t=k\Rightarrow t=k/2$~; $H=(k/2,k/2)=(\pm3\sqrt{2}/2,\pm3\sqrt{2}/2)$.

%────────────────────────────────────────────────────────────────
\noindent\phantomsection
{\color{BO}\bfseries\small D\'efi~4~\textendash\ \textit{Cercle inscrit dans un triangle}}
\par\smallskip

$A(0,4)$, $B(-3,0)$, $C(3,0)$.
$a=BC=6$~; $b=CA=\sqrt{9+16}=5$~; $c=AB=5$.
$I=\frac{aA+bB+cC}{a+b+c}=\frac{6(0,4)+5(-3,0)+5(3,0)}{16}=\frac{(0,24)}{16}=(0,\frac{3}{2})$.
$r=d(I,BC)=|y_I-0|=\frac{3}{2}$.
V\'erif.~: Aire $=\frac{1}{2}\cdot6\cdot4=12$~; $\frac{s}{1}=\frac{a+b+c}{2}=8$~;
$r=\frac{\text{Aire}}{s}=\frac{12}{8}=\frac{3}{2}$.\checkmark
Cercle inscrit~: $x^2+(y-\frac{3}{2})^2=\frac{9}{4}$.

%────────────────────────────────────────────────────────────────
\noindent\phantomsection
{\color{BO}\bfseries\small D\'efi~5~\textendash\ \textit{Puissance d'un point par rapport \`a un cercle}}
\par\smallskip

$P(4,3)$~; $\mathcal{C}$~: $x^2+y^2=25$.
$\pi(P)=4^2+3^2-25=25-25=0$~: $P$ est \textbf{sur} $\mathcal{C}$.\checkmark
$\pi(P)=0$~: cas limite (ni int\'erieur ni ext\'erieur).
$\pi(P)=x_P^2+y_P^2-r^2$ car $P\Omega^2=(x_P-0)^2+(y_P-0)^2=x_P^2+y_P^2$.\hfill$\square$
Axe radical~: $\pi_1(M)=\pi_2(M)\Leftrightarrow x^2+y^2+D_1x+E_1y+F_1=x^2+y^2+D_2x+E_2y+F_2$
$\Leftrightarrow(D_1-D_2)x+(E_1-E_2)y+(F_1-F_2)=0$.\hfill$\square$

%────────────────────────────────────────────────────────────────
\noindent\phantomsection
{\color{BO}\bfseries\small D\'efi~6~\textendash\ \textit{Th\'eor\`eme de la corde commune}}
\par\smallskip

$\mathcal{C}_1$~: $x^2+y^2=25$~; $\mathcal{C}_2$~: $(x-3)^2+y^2=16$.
Soustraction~: $x^2-(x-3)^2=9\Rightarrow6x-9=9\Rightarrow x=3$.
En $x=3$~: $y^2=25-9=16$~; $A=(3,4)$, $B=(3,-4)$.
Corde $[AB]$~: $AB=8$ (verticale, $x=3$).
Droite des centres~: $y=0$.
$x=3\perp y=0$.\checkmark
Angle entre corde et droite des centres~: $90^{\circ}$.

%────────────────────────────────────────────────────────────────
\noindent\phantomsection
{\color{BO}\bfseries\small D\'efi~7~\textendash\ \textit{\logiq\ Logique des droites de faisceau}}
\par\smallskip

\textbf{(1)} $\forall P\in\R^2,\;\exists$ une infinit\'e de droites passant par $P$
(une pour chaque direction, param\'etr\'ee par la pente $m\in\R\cup\{\infty\}$).

\textbf{(2)} Deux droites distinctes $d_1\ne d_2$~: si elles ont un point commun $Q$,
elles se croisent en $Q$. Supposons $Q$ et $Q'$ deux points communs distincts.
Alors $Q$ et $Q'$ d\'efinissent une droite unique~; $d_1$ et $d_2$ passent toutes les deux
par $Q$ et $Q'$~; elles co\"incident. Contradiction.\hfill$\square$

\textbf{(3)} Deux droites distinctes ne se coupent pas $\Rightarrow$ elles n'ont aucun point commun
$\Rightarrow$ elles sont parall\`eles. Raisonnement~: \textbf{contrapos\'ee} de (2).\hfill$\square$

%────────────────────────────────────────────────────────────────
\noindent\phantomsection
{\color{BO}\bfseries\small D\'efi~8~\textendash\ \textit{\algo\ G\'eom\'etrie analytique compl\`ete \computer}}
\par\smallskip

\begin{lstlisting}[language=Python]
import numpy as np, matplotlib.pyplot as plt, math

def analyse_triangle(A, B, C):
    G = ((A[0]+B[0]+C[0])/3, (A[1]+B[1]+C[1])/3)
    # Circocentre: intersection de deux mediataires
    mAB = [(A[0]+B[0])/2, (A[1]+B[1])/2]
    mBC = [(B[0]+C[0])/2, (B[1]+C[1])/2]
    dAB = [B[0]-A[0], B[1]-A[1]]
    dBC = [C[0]-B[0], C[1]-B[1]]
    M = np.array([[dAB[0],-dBC[0]],[dAB[1],-dBC[1]]], float)
    r = np.array([mBC[0]-mAB[0], mBC[1]-mAB[1]], float)
    t = np.linalg.solve(M, r)[0]
    Op = (mAB[0]+t*dAB[0], mAB[1]+t*dAB[1])
    R = math.hypot(Op[0]-A[0], Op[1]-A[1])
    # Orthocentre: H = A+B+C - 2*Op (relation d'Euler pour le circocentre=O')
    # En fait: H = A + B + C - 2*O' seulement si O est l'origine
    # Utiliser H = 3G - 2O'
    H = (3*G[0]-2*Op[0], 3*G[1]-2*Op[1])
    print(f"G={G}, O'={Op}, H={H}, R={R:.4f}")
    # Verification droite d'Euler
    aligned = abs((G[1]-Op[1])*(H[0]-Op[0]) - (H[1]-Op[1])*(G[0]-Op[0])) < 1e-9
    print(f"Droite d'Euler: {aligned}")
    # Trace
    theta = np.linspace(0, 2*np.pi, 300)
    plt.figure()
    tri = plt.Polygon([A,B,C], fill=False, edgecolor='b')
    plt.gca().add_patch(tri)
    plt.plot(*zip(A,B,C,A), 'b-')
    plt.plot([G[0],Op[0],H[0]],[G[1],Op[1],H[1]],'r--o')
    plt.plot(Op[0]+R*np.cos(theta), Op[1]+R*np.sin(theta),'g-')
    plt.axis('equal'); plt.grid(True); plt.show()

analyse_triangle((0,4),(-3,0),(3,0))
\end{lstlisting}

\needspace{3\baselineskip}
\needspace{3\baselineskip}
\needspace{3\baselineskip}
\needspace{3\baselineskip}
